\chapter{Hauptkapitel 2}
\label{sec: Hauptkapitel 2}

Dieser Abschnitt befasst sich insbesondere mit den eigenen Beiträgen zur Arbeit und kann nach Bedarf beispielsweise wiederum untergliedert sein in folgende Punkte:
\begin{itemize}
	\item Ist-Analysen
	\item Ergebnisse von Patentrecherchen
	\item Beschreibung des Standes der Technik
	\begin{itemize}
		\item Beschreibung von Arbeiten, die in direktem Zusammenhang stehen
		\item Informationen über bestehende Lösungen zu Teilen der Aufgabenstellung
		\item Diskussion des Bezuges zur eigenen Arbeit
		\item Darstellung der Vertrautheit mit dem Umfeld der eigenen Arbeit
		\item Eher faktenartige Darstellung
		\item Keine reinen Aufzählungen!
	\end{itemize}
	\item Erstellung eines Pflichtenheftes für den Sollzustand
	\item Berechnungen
	\item Tabellen und Diagramme
	\item Konzeptentwicklungen
	\item Prinzipskizzen und/oder Konstruktionszeichnungen
	\item Prototypentwicklung
	\item Problemlösungen und Versuchsergebnisse (insbesondere die eigenen Beiträge zur Problemlösung und die Ergebnisse so ausführlich wie notwendig):
	\begin{itemize}
		\item Analysen der Ausgangssituation und Festlegung der Herangehensweise an die Aufgabenstellung
		\item Entwicklung des Konzeptes und des grundsätzlichen Lösungsansatzes
		\item Detaillierung der Lösungserarbeitung mit allen Berechnungen, Konstruktionen, Programmschritten, Tabellen und Diagrammen \usw
		\item Ergebnisse, Erkenntnisse und Neuheiten
	\end{itemize}
	\item Bewertungen \usw
\end{itemize}	
Die originären Ideen, Innovationen sowie technologischen Neu- und Weiterentwicklungen sind entsprechend ausführlich zu behandeln und darzustellen.

\subsubsection*{Diskussion}
\begin{itemize}
	\item Diskussion der Ansätze, Vorgehensweise (Pro und Kontra)
	\item Ergebnisse, Erkenntnisse und deren Diskussion im Vergleich zu den Zielen
	\item Erklärung der Abweichungen
	\item Eigene Meinung zum Ergebnis
\end{itemize}
